\section{Introduction}
\label{sec:introduction}

Skyrmions provide a canonical setting in which topology, nonlinear field
equations, collective coordinates, and curved geometry meet
\cite{skyrme1961,manton2004}.  Rigorous analysis is not new here.  McLeod and
Troy proved existence and uniqueness for the spherically symmetric flat-space
Skyrme problem \cite{mcleodtroy1991}.  Creek, Donninger, Schlag, and Snelson
later gave a rational-arithmetic, mildly computer-assisted existence proof and
proved co-rotational linear stability on the whole space
\cite{creek2016}.  The present work is therefore not the first
computer-assisted existence or stability proof for a Skyrmion.

The nearest model variants are also established.  Static and gravitating
Skyrmions in de Sitter space, spinning collective-coordinate corrections, and
solitons confined by a mirror or cavity have been investigated numerically
\cite{brihaye2006,hata2010,giacomini2018}.  Nguyen and Nguyen explicitly survey
finite-radius boundary conditions for the original model
\cite{nguyen2013}.  These sources show that positive curvature, finite radius,
and validated numerics are not individually novel.  The narrower gap addressed
here is their conjunction for one massive, fixed-de-Sitter, unit-charge
Dirichlet hedgehog, including an augmented finite-boundary proof and certified
derived observables.

\begin{table}[ht]
  \centering
  \small
  \begin{tabular}{@{}p{0.20\linewidth}p{0.31\linewidth}p{0.43\linewidth}@{}}
    \toprule
    \raggedright Work & \raggedright Model &
    \raggedright Rigorous result or distinction \tabularnewline
    \midrule
    \raggedright McLeod--Troy \cite{mcleodtroy1991} &
    \raggedright Flat, whole-space spherical Skyrmion &
    \raggedright Existence and uniqueness. \tabularnewline
    \raggedright Creek et al.\ \cite{creek2016} &
    \raggedright Flat, whole-space co-rotational Skyrmion &
    \raggedright Rational computer-assisted existence and linear stability.
    \tabularnewline
    \raggedright Nguyen--Nguyen \cite{nguyen2013} &
    \raggedright Finite-radius boundary survey &
    \raggedright Numerical/global boundary analysis, not an interval
    certificate for the present massive curved problem. \tabularnewline
    \raggedright This work &
    \raggedright Massive fixed-de-Sitter hedgehog on a finite Dirichlet ball &
    \raggedright Local uniqueness in an augmented graph-norm ball,
    fixed-boundary radial gap, analytic local parameter branch, and
    authenticated boundary-induced inertia-transform tail. \tabularnewline
    \bottomrule
  \end{tabular}
  \caption{Closest rigorous and finite-radius comparisons.}
  \label{tab:prior-comparison}
\end{table}

This paper validates one deliberately narrow boundary-value problem.  We fix
the de Sitter metric rather than solving the Einstein equations, impose a
spherical Dirichlet wall strictly inside the cosmological horizon, and choose
one dimensionless parameter triple.  The central result is local: an
exact-rational Newton certificate proves a unique solution inside a stated
neighborhood of a rational spline.  The certificate also controls enough of
the solution to prove strict monotonicity, a negative wall slope, and positive
finite collective inertia.  Invertibility of the certified derivative then
gives a qualitative real-analytic solution branch on an open neighborhood of
the three parameters; the paper does not assign that neighborhood a numerical
width.

Three features make the validation more than a certified shooting zero.
First, the origin is regular singular in the usual radial variables.  We
replace it by a parameter-dependent Volterra chart and validate the regular
slope family and its first two sensitivities.  Second, the unknown shooting
slope is included in an augmented boundary map.  A coercive Dirichlet Jacobi
operator and a nonzero scalar Schur complement provide its inverse, while
row-local Green estimates and a center-equation Hessian identity close the
nonlinear tube.  Third, two consequences are authenticated against the same
profile certificate: a radial fluctuation gap and the exact large-frequency
tail of a finite optical transform of the rotor-inertia density.

The spectral result also makes the modeling boundary visible.  The optical
inertia weight has a quadratic zero at the Dirichlet wall, so boundary-aware
integration by parts gives an oscillatory $p^{-3}$ tail with nonzero
coefficient.  General endpoint expansions for finite Fourier and Hankel
transforms are classical \cite{wang2018,okadayamane2024}; our contribution is
not that regularity controls decay.  It is the certified endpoint order,
coefficient, and global derivative envelope for the exact nonlinear solution.

Computer-assisted contraction and radii-polynomial proofs for nonlinear
boundary-value problems are also established methodology
\cite{vandenberg2018,liu2021}.  Relative to Creek et al., the technical changes
are the finite regular-origin-to-Dirichlet matching problem, the unknown slope
as an augmented variable, a massive curved coefficient set, and authenticated
wall, inertia, radial-gap, and inertia-tail outputs.  The proof-relevant
implementation is a substantial exact-rational code base, inventoried in
\cref{sec:reproducibility}; it is not described as a small checker.  All
numerical values used in theorem statements are outward enclosures derived from
archived rational data; converged floating-point plots are used only for
orientation.

The paper is organized as follows.  \Cref{sec:model} defines the fixed
background, hedgehog problem, and central theorem.  \Cref{sec:validation}
develops the computer-assisted proof.  \Cref{sec:consequences} proves the
radial gap and spectral tail.  \Cref{sec:reproducibility} identifies the
source-pinned validation surface, archives, and reproduction path.
\Cref{sec:discussion} states the
mathematical role of the Dirichlet boundary and the limits of the result.
